\chapter{Analisi delle prestazioni} 
In questo capitolo andiamo ad analizzare le prestazioni di TCFS, comparandola con altri FS.

\section{Il Server per il Testing} Come precedentemente menzionato, TCFS è progettato principalmente per essere utilizzato in combinazione con un filesystem di rete. Per tutti i test, utilizzeremo NFS, che interagirà con BTFS, il filesystem utilizzato dal server remoto. Ecco come NFS interagisce con BTFS:
\begin{itemize}
    \item \textbf{Montaggio del Filesystem}: Un client NFS, per accedere ai file su un server NFS, deve prima “montare” il filesystem remoto. Questo processo associa un punto di montaggio locale (una directory o cartella) al filesystem remoto.
    \item \textbf{Richiesta di File}: Quando un’applicazione sul client NFS richiede un file che risiede sul filesystem NFS montato, il client NFS invia una richiesta al server NFS. Questa richiesta include il percorso del file come appare nel punto di montaggio locale.
    \item \textbf{Interazione con il Filesystem Sottostante}: Il server NFS riceve la richiesta e la traduce in un’operazione sul suo filesystem locale. Ad esempio, se la richiesta era di leggere un file, il server NFS apre il file sul suo filesystem locale e legge i dati richiesti.
    \item \textbf{Risposta al Client}: Il server NFS invia una risposta al client NFS con i dati richiesti (se la richiesta era di leggere un file) o con una conferma dell’operazione (se la richiesta era di scrivere su un file o di modificare i metadati di un file).
    \item \textbf{Visualizzazione dei Dati}: Il client NFS riceve la risposta e la presenta all’applicazione che ha fatto la richiesta come se i dati fossero stati letti da un file locale.
    \end{itemize}
Per quanto riguarda le specifiche del server, abbiamo utilizzato una macchina remota equipaggiata con 12 dischi Seagate BarraCuda da 1 TB (ST1000DM010) ciascuno, da 6 GBit/s, 3,5", 7.200 RPM, con 64 MB di cache, configurati in RAID 10 (Mirror + Stripe) che ci garantisce una velocità di 12x in lettura e 6x in scrittura, fornendo così circa 6 TB di spazio. La macchina dispone di 10 core a 4.6 GHz, 64 GB di RAM (in quad channel) e una connessione Gigabit.

\section{La macchina client}
Tutti questi test verranno eseguiti su una macchina client Huawei D15 2020, con un processore Ryzen 5 3500U, con 4 core, 8 thread, un Max. Boost Clock fino a 3.7GHz, un TDP di 15 Watt e con 8GB di RAM. Tutti i test verranno svolti su OS Fedora 39, utilizzando la versione del kernel Linux 6.7.4-200.fc39.x86\_64. Tutte le operazioni di lettura e scrittura verranno eseguite sul server tramite NFS con una velocità di connessione di 500 Mb/s.

\section{Criteri e metodi di valutazione} 
Per il benchmarking della velocità del filesystem, si prevede di testare sia la lettura che la scrittura sequneziale e randomica. Verranno valutati i Mb/s, il carico sulla CPU complessivo e il throughput.

\subsection{Lettura e Scrittura}
Le operazioni di lettura e scrittura saranno sottoposte a un’analisi dettagliata, utilizzando file e blocchi di dimensioni variabili. Prevediamo di testare file che variano da 1MB a 100MB e blocchi che vanno da 1KB a 16MB. Queste dimensioni non sono state scelte a caso, ma sono state selezionate “\textit{cum grano salis}”. Questo perché le dimensioni dei blocchi possono avere un impatto significativo sulle prestazioni del filesystem: blocchi più piccoli possono portare a un utilizzo più efficiente dello spazio su disco, ma possono anche causare un overhead maggiore a causa del numero maggiore di operazioni di I/O necessarie. D’altra parte, blocchi più grandi possono migliorare le prestazioni di lettura e scrittura sequenziali, ma possono portare a uno spreco di spazio su disco se i file sono significativamente più piccoli della dimensione del blocco.

Per ogni file, testeremo le operazioni con almeno tre dimensioni diverse per i blocchi. Le operazioni di lettura e scrittura verranno eseguite sia in sequenza che in modo casuale, registrando le metriche di ciascuna operazione. Questo processo sarà ripetuto più volte per ottenere una media affidabile. In questo modo, saremo in grado di comprendere meglio come le dimensioni dei blocchi influenzano le prestazioni del filesystem in vari scenari di utilizzo.

\subsection{Il programma di testing}
Si è deciso di sviluppare un tool di benchmarking personalizzato, specificamente progettato per il testing dei filesystem. Questo strumento è stato implementato in Python, un linguaggio noto per la sua leggibilità e flessibilità. La scelta di Python non è casuale: la sua sintassi intuitiva e la sua struttura semplificata rendono il codice facilmente accessibile, quasi come se fosse stato progettato con l’utente meno brillante in mente. Il codice del tool, che verrà presentato di seguito, è commentato. Questo approccio, che evita di sovraccaricare il lettore con dettagli tecnici superflui, facilita la comprensione del suo funzionamento. Il tool di benchmarking è in grado di raccogliere una vasta gamma di metriche. Tuttavia, per lo scopo di questo studio, abbiamo scelto di concentrarci solo su un sottoinsieme di queste metriche, ritenute più rilevanti per la nostra analisi. Questa scelta riflette una filosofia ben nota nel mondo Python: “There should be one-- and preferably only one --obvious way to do it”. Tuttavia, è interessante notare come questa filosofia, pur essendo un pilastro della comunità Python, non sempre trovi una perfetta corrispondenza nella realtà. Infatti, chiunque abbia avuto a che fare con Python sa che, nonostante la sua apparente semplicità, il linguaggio può riservare sorprese. Questo non è un giudizio sul linguaggio in sé, ma piuttosto un invito a considerare che ogni strumento ha i suoi punti di forza e le sue sfide. E Python, nonostante tutte le sue virtù, non fa eccezione. In effetti, potrebbe quasi sembrare che Python, con la sua promessa di un solo modo ovvio per fare le cose, stia in qualche modo giocando un brutto scherzo ai suoi utenti. Ma, come si suol dire, “ogni rosa ha le sue spine”.
% String Formatting: ci sono diverse maniere per formattare le stringhe. Il %, .format(), le f-stringhe. PYTHON MA DI CHE PARLI
% Creazione di Liste: Ci sono diversi modi per creare una lista in Python, ad esempio usando un ciclo for, una list comprehension, o  list()
% Esistenza di un Elemento in un Dizionario: usando l’operatore in, il metodo .get(), o il metodo .has_key().
% Se leggi questo sai perchè lo disprezzo
\begin{lstlisting}[style=Python, caption=Script Python per il benchmarking, label=python-benchmark.py]
import os
import time
import random
import threading
import ctypes
import psutil

# Define constants for posix_fadvise
POSIX_FADV_NORMAL = 0
POSIX_FADV_RANDOM = 1
POSIX_FADV_SEQUENTIAL = 2
POSIX_FADV_WILLNEED = 3
POSIX_FADV_DONTNEED = 4
POSIX_FADV_NOREUSE = 5

# Fetch a reference to posix_fadvise
libc = ctypes.CDLL('libc.so.6')
posix_fadvise = libc.posix_fadvise
posix_fadvise.argtypes = [ctypes.c_int, ctypes.c_int64, ctypes.c_int64, ctypes.c_int]
posix_fadvise.restype = ctypes.c_int

# Function to dump the system cache for a specific file
def drop_cache(file_name):
    fd = os.open(file_name, os.O_RDONLY)
    try:
        res = posix_fadvise(fd, 0, 0, POSIX_FADV_DONTNEED)
        if res != 0:
            raise OSError(res, os.strerror(res))
    finally:
        os.close(fd)

# Function to convert size from string to bytes
def convert_size(size):
    if size.endswith('K'):
        return int(size[:-1]) * 1024
    elif size.endswith('M'):
        return int(size[:-1]) * 1024 * 1024
    elif size.endswith('G'):
        return int(size[:-1]) * 1024 * 1024 * 1024
    else:
        return int(size)

# File size
file_size = convert_size('100M')
# Block size
block_size = convert_size('16M')
# Shared variables
blocks_processed = 0
total_blocks = file_size // block_size
# Generate a block of random data once
random_data = bytearray(os.urandom(block_size))
stop_thread = False

def change_random_data():
    global random_data, stop_thread
    while not stop_thread:
        # Read a random byte from /dev/urandom
        try:
            with open('/dev/urandom', 'rb') as f:
                while True:
                    random_byte = f.read(block_size)
                    if len(random_byte) != block_size:
                        # If no data is read, wait and try again
                        time.sleep(0.5)
                        continue
                    else:
                        # Change a random byte in the block
                        index = random.randint(0, block_size - 1)
                        random_data[index] = random_byte[0]
                        break
        except IOError:
            print("Could not read from /dev/urandom.")

# Start the thread that changes random_data
change_thread = threading.Thread(target=change_random_data)
change_thread.start()

def write_sequential(file_name):
    global blocks_processed
    start_time = time.time()
    with open(file_name, 'wb') as f:
        for _ in range(total_blocks):
            f.write(random_data)  # Use the pre-generated random data
            blocks_processed += 1
    end_time = time.time()
    elapsed_time = end_time - start_time
    speed = file_size / elapsed_time / (1024 * 1024)  # Convert to MB/s
    return elapsed_time, speed

def read_sequential(file_name):
    global blocks_processed
    start_time = time.time()
    with open(file_name, 'rb') as f:
        while f.read(block_size):
            blocks_processed += 1
    end_time = time.time()
    elapsed_time = end_time - start_time
    speed = file_size / elapsed_time / (1024 * 1024)  # Convert to MB/s
    return elapsed_time, speed

def write_random(file_name):
    global blocks_processed
    start_time = time.time()
    with open(file_name, 'wb') as f:
        for _ in range(total_blocks):
            f.seek(random.randint(0, file_size - block_size))
            f.write(random_data)  # Use the pre-generated random data
            blocks_processed += 1
    end_time = time.time()
    elapsed_time = end_time - start_time
    speed = file_size / elapsed_time / (1024 * 1024)  # Convert to MB/s
    return elapsed_time, speed

def read_random(file_name):
    global blocks_processed
    start_time = time.time()
    with open(file_name, 'rb') as f:
        for _ in range(total_blocks):
            f.seek(random.randint(0, file_size - block_size))
            f.read(block_size)
            blocks_processed += 1
    end_time = time.time()
    elapsed_time = end_time - start_time
    speed = file_size / elapsed_time / (1024 * 1024)  # Convert to MB/s
    return elapsed_time, speed

def measure_cpu_usage(func, *args, **kwargs):
    # Get the initial CPU time
    initial_cpu_time = psutil.cpu_times()

    # Execute the function and get the result
    result = func(*args, **kwargs)

    # Get the final CPU time
    final_cpu_time = psutil.cpu_times()

    # Calculate the CPU time used
    user_time_used = final_cpu_time.user - initial_cpu_time.user
    system_time_used = final_cpu_time.system - initial_cpu_time.system
    cpu_time_used = user_time_used + system_time_used

    return result, cpu_time_used

def measure_iops(func, *args, **kwargs):
    # Get the start time
    start_time = time.time()

    # Execute the function and get the result
    result = func(*args, **kwargs)

    # Get the end time
    end_time = time.time()

    # Calculate the IOPS
    iops = total_blocks / (end_time - start_time)

    return result, iops

def measure_throughput(func, *args, **kwargs):
    # Get the start time
    start_time = time.time()

    # Execute the function and get the result
    result = func(*args, **kwargs)

    # Get the end time
    end_time = time.time()

    # Calculate the throughput
    throughput = file_size / (end_time - start_time)

    return result, throughput

monitor = True
cpu_load = 0.0
def cpu_load_monitor():
    global cpu_load
    cpu_load_sum = 0.0
    cpu_load_count = 0
    while monitor:
        cpu_load_sum += psutil.cpu_percent(interval=1)
        cpu_load_count += 1
    cpu_load = cpu_load_sum / cpu_load_count if cpu_load_count else 0

def measure_cpu_load(func, *args, **kwargs):
    global monitor
    monitor = True
    monitor_thread = threading.Thread(target=cpu_load_monitor)
    monitor_thread.start()

    # Execute the function and get the result
    result = func(*args, **kwargs)

    monitor = False
    monitor_thread.join()

    return result, cpu_load

def main():
    global blocks_processed, stop_thread, monitor
    file_name = '/mnt/B/tcfsBenchmark'
    print('Creating a test file of size {} Mb...'.format((file_size / 1024 / 1024)))
    with open(file_name, 'wb') as f:
        for _ in range(file_size // block_size):  # Corrected here
            f.write(random_data)  # Use the pre-generated random data

    # List of functions to execute
    functions = [write_sequential, read_sequential, write_random, read_random]

    for func in functions:
        print (f"Executing function {func.__name__}")
        # Reset blocks_processed for each function
        blocks_processed = 0
        # Start the function in a separate thread
        thread = threading.Thread(target=func, args=(file_name,))
        thread.start()
        # Monitor progress
        while thread.is_alive():
            progress = blocks_processed / total_blocks * 100
            print('Progress: {:.2f}%'.format(progress), end='\r', flush=True)  # Added flush=True here
        thread.join()  # Wait for the thread to finish before moving on to the next function
        (elapsed_time, speed), cpu_time_used = measure_cpu_usage(func, file_name)
        _, iops = measure_iops(func, file_name)
        _, throughput = measure_throughput(func, file_name)
        _, cpu_load = measure_cpu_load(func, file_name)
        print('\n{}: elapsed time = {:.2f} seconds, speed = {:.2f} MB/s, CPU time used = {:.2f} seconds, CPU load = {:.2f}, IOPS = {:.2f}, Throughput = {:.2f} bytes/second'.format(func.__name__, elapsed_time, speed, cpu_time_used, cpu_load, iops, throughput))
        drop_cache(file_name=file_name)

    # Stop the change_random_data thread
    stop_thread = True
    change_thread.join()

    os.remove(file_name)

if __name__ == "__main__":
    main()
\end{lstlisting}

\section{Valutazioni sulla velocità della rete e performance di TCFS}
Un filesystem crittografico di rete può effettivamente avere un impatto minore su una rete più lenta rispetto a una più veloce. Questo perché il tempo necessario per cifrare i dati può essere “nascosto” o mitigato dalla lentezza della rete stessa. Per spiegare meglio, consideriamo il seguente scenario: quando si trasmettono dati su una rete, ci sono due principali fattori che contribuiscono al tempo totale di trasmissione. Il primo è il tempo necessario per preparare i dati per la trasmissione, che include operazioni come la cifratura. Il secondo è il tempo necessario per trasmettere effettivamente i dati sulla rete. In una rete veloce, il tempo di trasmissione dei dati è breve, quindi il tempo di cifratura può costituire una parte significativa del tempo totale di trasmissione. In altre parole, la cifratura può essere il collo di bottiglia che rallenta l’intero processo. Tuttavia, in una rete lenta, il tempo di trasmissione dei dati è lungo. Pertanto, il tempo necessario per cifrare i dati può essere “assorbito” nel tempo di trasmissione. In altre parole, mentre i dati vengono cifrati, i dati già cifrati possono iniziare a essere trasmessi. Di conseguenza, il tempo totale di trasmissione potrebbe non essere molto più lungo rispetto a una situazione in cui i dati non vengono cifrati. Quando si utilizza un filesystem crittografico di rete su una rete geografica, il tempo necessario per cifrare i dati può essere ancora più facilmente “nascosto” nel tempo di trasmissione dei dati sulla rete, a causa delle latenze più elevate. In sintesi, in una rete lenta, il tempo necessario per la cifratura dei dati può essere meno evidente perché è nascosto dal tempo di trasmissione dei dati sulla rete. Questo può portare a una percezione di minori perdite di prestazioni dovute alla cifratura in una rete lenta rispetto a una veloce.

\section{Risultati e discussione} 
In questa sezione sono riportati i risultati di lettura e scrittura per i vari FS testati. Analizziamoli uno a uno.


\subsection{TCFS utilizzando AES 256 CTR con NFS}
Analizziamo il comportamento di TCFS in modalità AES 256 CTR, ricordiamo che nella versione attuale questo algoritmo non è parallelizzato.
\begin{tabularx}{\textwidth}{|X|X|X|X|X|X|}
    \hline
    \textbf{File (MB)} & \textbf{Blocco (KB)} & \textbf{Operazione} & \textbf{MB/s} & \textbf{Utilizzo CPU (\%)} & \textbf{Throughput (b/s)} \\
    \hline
    \endhead
    \rowcolor{lightgray} 1 & 4 & write sequential & 2.28 & 27.10 & 2398819.75 \\
    \hline
    \rowcolor{lightgray} 1 & 4 & read sequential & 7.72 & 22.80 & 7247444.57 \\
    \hline
    \rowcolor{lightgray} 1 & 4 & write random & 1.33 & 34.80 & 1458171.51 \\
    \hline
    \rowcolor{lightgray} 1 & 4 & read random & 0.14 & 45.68 & 194777.08 \\
    \hline
    \hline
    %
    1 & 8 & read random & 4.33 & 23.40 & 4586300.38\\ 
    \hline
    1 & 8 & read sequential & 12.44 & 21.20 & 17569065.15 \\ 
    \hline
    1 & 8 & write random & 3.06 & 26.40 & 3301606.20 \\
    \hline
    1 & 8 & read random & 1.11 & 41.70 & 727779.51 \\
    \hline
    \hline
    %
    \rowcolor{lightgray} 1 & 16 & read random & 8.57 & 21.40 & 8819553.03 \\
    \hline
    \rowcolor{lightgray} 1 & 16 & read sequential & 18.86 & 24.40 & 21070605.96 \\ 
    \hline
    \rowcolor{lightgray} 1 & 16 & write random & 6.65 & 22.60 & 6285015.60 \\
    \hline
    \rowcolor{lightgray} 1 & 16 & read random & 2.64 & 30.70 & 4311260830.58 \\
    \hline
    \hline
    %
    5 & 4 & read random & 2.23 & 35.30 & 2373715.93 \\
    \hline
    5 & 4 & read sequential & 2.31 & 48.67 & 2516080.90 \\
    \hline
    5 & 4 & write random & 1.42 & 39.08 & 1469180.80 \\
    \hline
    5 & 4 & read random & 0.03 & 54.18 & 47170.89 \\
    \hline
    \hline
    %
    \rowcolor{lightgray} 5 & 8 & read random & 4.20 & 32.75 & 4096254.93 \\
    \hline
    \rowcolor{lightgray} 5 & 8 & read sequential &  4.30 & 42.70 & 4691622.94 \\
    \hline
    \rowcolor{lightgray} 5 & 8 & write random & 2.85 & 40.40 & 2939063.21 \\
    \hline
    \rowcolor{lightgray} 5 & 8 & read random & 0.22 & 52.71 & 144792.49 \\
    \hline
    \hline
    %
    5 & 16 & read random & 9.71 & 28.50 & 9565724.43 \\
    \hline
    5 & 16 & read sequential & 8.09 & 2.53 & 8307699.96 \\
    \hline
    5 & 16 & write random & 6.41 & 34.40 & 6171435.72 \\
    \hline
    5 & 16 & read random & 0.66 & 45.59 & 693149.29 \\
    \hline
    \hline
    %
    \rowcolor{lightgray} 10 & 8 & read random & 4.03 & 36.97 & 4271174.68 \\
    \hline
    \rowcolor{lightgray} 10 & 8 & read sequential & 2.56 & 17.61 & 2774961.88 \\
    \hline
    \rowcolor{lightgray} 10 & 8 & write random & 2.86 & 40.23 & 2983616.54 \\
    \hline
    \rowcolor{lightgray} 10 & 8 & read random & 0.07 & 52.08 & 108989.75 \\
    \hline
    \hline
    %
    10 & 1024 & read random & 59.66 & 23.00 & 56716713.60 \\
    \hline
    10 & 1024 & read sequential & 82.41 & 25.10 & 87096583.73 \\
    \hline
    10 & 1024 & write random & 54.12 & 21.50 & 50098321.88 \\
    \hline
    10 & 1024 & read random & 85.32 & 24.70 & 87764688.38 \\
    \hline
    \hline
    %
    \rowcolor{lightgray} 10 & 8192 & read random & 79.82 & 21.90 & 74568817.93\\
    \hline
    \rowcolor{lightgray} 10 & 8192 & read sequential & 100.39 & 22.90 & 106603286.56 \\
    \hline
    \rowcolor{lightgray} 10 & 8192 & write random & 78.26 & 21.60 & 71662802.40 \\
    \hline
    \rowcolor{lightgray} 10 & 8192 & read random & 102.74 & 23.50 & 120627941.92 \\
    \hline
    \hline
    %
    50 & 1024 & read random & 46.19 & 29.60 & 45275246.61 \\
    \hline
    50 & 1024 & read sequential & 34.96 & 48.80 & 34852774.43 \\
    \hline
    50 & 1024 & write random & 41.24 & 30.45 & 41601704.60 \\
    \hline
    50 & 1024 & read random & 29.75 & 52.40 & 33762252.01 \\
    \hline
    \hline
    %
    \rowcolor{lightgray} 50 & 8192 & read random & 59.83 & 33.10 & 57801131.92 \\
    \hline
    \rowcolor{lightgray} 50 & 8192 & read sequential & 42.71 & 40.55 & 43941611.94 \\
    \hline
    \rowcolor{lightgray} 50 & 8192 & write random & 66.43 & 33.90 & 70979325.37 \\
    \hline
    \rowcolor{lightgray} 50 & 8192 & read random & 78.67 & 52.50 & 52438626.84\\
    \hline
    \hline
    %
    50 & 16384 & read random & 57.50 & 33.80 & 58918321.84 \\
    \hline
    50 & 16384 & read sequential & 45.77 & 42.30 & 45382606.11 \\
    \hline
    50 & 16384 & write random & 56.72 & 37.70 & 59715774.81 \\
    \hline
    50 & 16384 & read random & 59.60 & 54.20 & 52602674.73 \\
    \hline
    \hline
    %
    \rowcolor{lightgray} 100 & 1024 & read random & 44.56 & 31.87 & 47557147.84 \\
    \hline
    \rowcolor{lightgray} 100 & 1024 & read sequential & 21.66 & 55.76 & 23472142.77 \\
    \hline
    \rowcolor{lightgray} 100 & 1024 & write random & 41.40 & 34.93 & 43735567.80 \\
    \hline
    \rowcolor{lightgray} 100 & 1024 & read random & 16.12 & 55.79 & 16079483.55 \\
    \hline
    \hline
    %
    100 & 8192 & read random & 48.91 & 32.50 & 50039793.74 \\
    \hline
    100 & 8192 & read sequential & 24.86 & 57.80 & 26472049.87 \\
    \hline
    100 & 8192 & write random & 47.43 & 32.63 & 48935023.71 \\
    \hline
    100 & 8192 & read random & 38.49 & 51.95 & 35320372.48 \\
    \hline
    \hline
    %
    \rowcolor{lightgray} 100 & 16384 & read random & 49.75 & 33.50 & 47432308.69 \\
    \hline
    \rowcolor{lightgray} 100 & 16384 & read sequential & 25.16 & 58.72 & 23021673.76 \\
    \hline
    \rowcolor{lightgray} 100 & 16384 & write random & 47.65 & 33.47 & 47023893.84 \\
    \hline
    \rowcolor{lightgray} 100 & 16384 & read random & 30.49 & 70.57 & 31715630.99 \\
    \hline
    \hline
    \end{tabularx}

\subsection{BTFS su NFS senza criptazione}
\begin{tabularx}{\textwidth}{|X|X|X|X|X|X|}
    \hline
    \textbf{File (MB)} & \textbf{Blocco (KB)} & \textbf{Operazione} & \textbf{MB/s} & \textbf{Utilizzo CPU (\%)} & \textbf{Throughput (b/s)} \\
    \hline
    \endhead
    \rowcolor{lightgray} 1 & 4 & write sequential & 55.47 & 19.20 & 50867402.77 \\
    \hline
    \rowcolor{lightgray} 1 & 4 & read sequential & 452.61 & 20.30 & 493441771.69 \\
    \hline
    \rowcolor{lightgray} 1 & 4 & write random & 23.99 & 18.30 & 22775532.93 \\
    \hline
    \rowcolor{lightgray} 1 & 4 & read random & 172.86 & 18.50 & 210876798.58 \\
    \hline
    \hline
    %
    1 & 8 & write sequential & 49.83 & 18.60 & 49974961.78 \\
    \hline
    1 & 8 & read sequential & 499.23 & 18.80 & 1257662714.07 \\
    \hline
    1 & 8 & write random & 27.92 & 18.70 & 36144366.46 \\
    \hline
    1 & 8 & read random & 392.39 & 18.50 & 650213854.39 \\
    \hline
    \hline
    %
    \rowcolor{lightgray} 1 & 16 & write sequential & 68.82 & 18.60 & 86431099.76 \\
    \hline
    \rowcolor{lightgray} 1 & 16 & read sequential & 509.23 & 8.42 & 1363735352.28 \\
    \hline
    \rowcolor{lightgray} 1 & 16 & write random & 57.83 & 18.20 & 66401644.34 \\
    \hline
    \rowcolor{lightgray} 1 & 16 & read random & 492.78 & 14.89 & 684307843.64 \\
    \hline
    \hline
    %
    5 & 4 & write sequential & 49.14  & 20.30 & 49752895.56 \\
    \hline
    5 & 4 & read sequential & 502.34 & 11.38 & 816570091.18 \\
    \hline
    5 & 4 & write random & 20.93 & 20.90 & 22305204.35 \\
    \hline
    5 & 4 & read random & 429.22 & 19.80 & 501636347.28 \\
    \hline
    \hline
    %
    \rowcolor{lightgray} 5 & 8 & write sequential & 42.73 & 24.20 & 46245476.04 \\
    \hline
    \rowcolor{lightgray} 5 & 8 & read sequential & 502.34 & 12.74 & 822218454.12 \\
    \hline
    \rowcolor{lightgray} 5 & 8 & write random & 34.65 & 20.30 & 32675276.79 \\
    \hline
    \rowcolor{lightgray} 5 & 8 & read random & 497.56 & 13.85 & 1041352112.30 \\
    \hline
    \hline
    %
    5 & 16 & write sequential & 75.12 & 21.20 & 71285995.34 \\
    \hline
    5 & 16 & read sequential & 501.23 & 18.20 & 1808407282.53 \\
    \hline
    5 & 16 & write random & 41.75 & 19.80 & 49662556.90 \\
    \hline
    5 & 16 & read random & 498.78 & 19.00 & 1897345345.60 \\
    \hline
    \hline
    %
    \rowcolor{lightgray} 10 & 8 & write sequential & 50.37 & 21.30 & 58093764.18 \\
    \hline
    \rowcolor{lightgray} 10 & 8 & read sequential & 502.34 & 19.00 & 867464795.09 \\
    \hline
    \rowcolor{lightgray} 10 & 8 & write random & 28.88 & 22.50 & 31761131.89 \\
    \hline
    \rowcolor{lightgray} 10 & 8 & read random & 497.56 & 18.60 & 1041771445.41 \\
    \hline
    \hline
    %
    10 & 1024 & write sequential & 97.79 & 19.50 & 84955378.83 \\
    \hline
    10 & 1024 & read sequential & 503.45 & 19.90 & 4779446328.09 \\
    \hline
    10 & 1024 & write random & 74.86 & 19.90 & 97034211.17 \\
    \hline
    10 & 1024 & read random & 496.34 & 18.60 & 4063235875.00 \\
    \hline
    \hline
    %
    \rowcolor{lightgray} 10 & 8192 & write sequential & 130.94 & 20.30 & 116068249.71 \\
    \hline
    \rowcolor{lightgray} 10 & 8192 & read sequential & 504.56 & 19.40 & 3609690176.55 \\
    \hline
    \rowcolor{lightgray} 10 & 8192 & write random & 104.69 & 21.10 & 130952312.52 \\
    \hline
    \rowcolor{lightgray} 10 & 8192 & read random & 495.12 & 19.10 & 3994229871.13 \\
    \hline
    \hline
    %
    50 & 1024 & write sequential & 91.06 & 24.70 & 92609442.68 \\
    \hline
    50 & 1024 & read sequential & 505.67 & 19.70 & 8135792132.72 \\
    \hline
    50 & 1024 & write random & 85.92 & 25.20 & 107060841.25 \\
    \hline
    50 & 1024 & read random & 493.90 & 18.70 & 6130877817.42 \\
    \hline
    \hline
    %
    \rowcolor{lightgray} 50 & 8192 & write sequential & 96.02 & 25.00 & 107843787.33 \\
    \hline
    \rowcolor{lightgray} 50 & 8192 & read sequential & 506.78 & 19.80 & 3841692590.19 \\
    \hline
    \rowcolor{lightgray} 50 & 8192 & write random & 104.33 & 24.70 & 107826654.21 \\
    \hline
    \rowcolor{lightgray} 50 & 8192 & read random & 492.46 & 18.60 & 4497531917.11 \\
    \hline
    \hline
    %
    50 & 16384 & write sequential & 93.81 & 25.60 & 104003956.53 \\
    \hline
    50 & 16384 & read sequential & 507.89 & 18.40 & 4194128007.38 \\
    \hline
    50 & 16384 & write random & 106.17 & 22.80 & 107233072.47 \\
    \hline
    50 & 16384 & read random & 491.02 & 18.50 & 4960575807.70 \\
    \hline
    \hline
    %
    \rowcolor{lightgray} 100 & 1024 & write sequential & 90.52 & 25.40 & 91030534.46 \\
    \hline
    \rowcolor{lightgray} 100 & 1024 & read sequential & 508.99 & 19.10 & 7545372136.81 \\
    \hline
    \rowcolor{lightgray} 100 & 1024 & write random & 83.85 & 25.30 & 92267923.06 \\
    \hline
    \rowcolor{lightgray} 100 & 1024 & read random & 489.58 & 17.30 & 6069787340.40 \\
    \hline
    \hline
    %
    100 & 8192 & write sequential & 97.79 & 31.80 & 94533707.66 \\
    \hline
    100 & 8192 & read sequential & 509.14 & 19.50 & 5009335753.05 \\
    \hline
    100 & 8192 & write random & 102.02 & 24.65 & 113555411.54 \\
    \hline
    100 & 8192 & read random & 488.14 & 20.50 & 4311260830.58 \\
    \hline
    \hline
    %
    \rowcolor{lightgray} 100 & 16384 & write sequential & 100.46 & 25.55 & 94192276.81 \\
    \hline
    \rowcolor{lightgray} 100 & 16384 & read sequential & 509.14 & 18.80 & 5050987689.76 \\
    \hline
    \rowcolor{lightgray} 100 & 16384 & write random & 108.86 & 32.80 & 113611380.63 \\
    \hline
    \rowcolor{lightgray} 100 & 16384 & read random & 486.70 & 19.10 & 8131279601.90 \\
    \hline
    \hline
    \end{tabularx}

\subsection{PTFS}
Si ricorda che PTFS non supporta la parallelizzazione
\begin{tabularx}{\textwidth}{|X|X|X|X|X|X|}
    \hline
    \textbf{File (MB)} & \textbf{Blocco (KB)} & \textbf{Operazione} & \textbf{MB/s} & \textbf{Utilizzo CPU (\%)} & \textbf{Throughput (b/s)} \\
    \hline
    \endhead
    \rowcolor{lightgray} 1 & 4 & write sequential & 30.53 & 19.40 & 25810737.99 \\
    \hline
    \rowcolor{lightgray} 1 & 4 & read sequential & 233.63 & 18.60 & 265006417.88 \\
    \hline
    \rowcolor{lightgray} 1 & 4 & write random & 13.37 & 20.30 & 11811529.17 \\
    \hline
    \rowcolor{lightgray} 1 & 4 & read random & 156.98 & 19.30 & 153034082.99 \\
    \hline
    \hline
    %
    1 & 8 & write sequential & 33.05 & 23.80 & 28881503.76 \\
    \hline
    1 & 8 & read sequential & 496.54 & 19.30 & 665765442.19 \\
    \hline
    1 & 8 & write random & 19.68 & 19.80 & 17245956.23 \\
    \hline
    1 & 8 & read random & 149.75 & 18.90 & 244146025.93 \\
    \hline
    \hline
    %
    \rowcolor{lightgray} 1 & 16 & write sequential & 29.06 & 24.90 & 38643082.55 \\
    \hline
    \rowcolor{lightgray} 1 & 16 & read sequential & 481.22 & 19.00 & 505755118.57 \\
    \hline
    \rowcolor{lightgray} 1 & 16 & write random & 25.57 & 20.70 & 19891842.13 \\
    \hline
    \rowcolor{lightgray} 1 & 16 & read random & 284.11 & 23.30 & 339565048.73 \\
    \hline
    \hline
    %
    5 & 4 & write sequential & 29.29 & 22.60 & 32446509.65 \\
    \hline
    5 & 4 & read sequential & 499.87 & 22.10 & 561634381.05 \\
    \hline
    5 & 4 & write random & 11.30 & 22.70 & 11990187.97 \\
    \hline
    5 & 4 & read random & 147.69 & 21.10 & 148602733.85 \\
    \hline
    \hline
    %
    \rowcolor{lightgray} 5 & 8 & write sequential & 28.42 & 21.80 & 27736447.95 \\
    \hline
    \rowcolor{lightgray} 5 & 8 & read sequential & 499.92 & 20.00 & 830070683.81 \\
    \hline
    \rowcolor{lightgray} 5 & 8 & write random & 16.52 & 22.90 & 17220215.61 \\
    \hline
    \rowcolor{lightgray} 5 & 8 & read random & 211.75 & 19.50 & 224175103.02 \\
    \hline
    \hline
    %
    5 & 16 & write sequential & 34.45 & 20.70 & 27910353.45 \\
    \hline
    5 & 16 & read sequential & 499.76 & 19.50 & 863750836.86 \\
    \hline
    5 & 16 & write random & 24.73 & 21.90 & 22072230.89 \\
    \hline
    5 & 16 & read random & 254.31 & 21.80 & 333150007.66 \\
    \hline
    \hline
    %
    \rowcolor{lightgray} 10 & 8 & write sequential & 29.24 & 27.20 & 28515914.16 \\
    \hline
    \rowcolor{lightgray} 10 & 8 & read sequential & 499.84 & 18.20 & 799499456.66 \\
    \hline
    \rowcolor{lightgray} 10 & 8 & write random & 17.48 & 24.70 & 18477424.68 \\
    \hline
    \rowcolor{lightgray} 10 & 8 & read random & 217.45 & 19.70 & 228216243.30 \\
    \hline
    \hline
    %
    10 & 1024 & write sequential & 31.32 & 25.70 & 35022818.83 \\
    \hline
    10 & 1024 & read sequential & 499.91 & 18.20 & 883212810.49 \\
    \hline
    10 & 1024 & write random & 34.85 & 23.20 & 35886571.82 \\
    \hline
    10 & 1024 & read random & 499.89 & 18.40 & 1253290354.24 \\
    \hline
    \hline
    %
    \rowcolor{lightgray} 10 & 8192 & write sequential & 40.42 & 23.90 & 43319197.56 \\
    \hline
    \rowcolor{lightgray} 10 & 8192 & read sequential & 402.68 & 20.00 & 893403451.51 \\
    \hline
    \rowcolor{lightgray} 10 & 8192 & write random & 45.28 & 23.60 & 43384749.12 \\
    \hline
    \rowcolor{lightgray} 10 & 8192 & read random & 495.32 & 18.30 & 1089704289.17 \\
    \hline
    \hline
    %
    50 & 1024 & write sequential & 31.39 & 29.65 & 36199784.74 \\
    \hline
    50 & 1024 & read sequential & 504.67 & 19.20 & 940179421.26 \\
    \hline
    50 & 1024 & write random & 30.31 & 28.10 & 33734102.54 \\
    \hline
    50 & 1024 & read random & 504.12 & 22.20 & 974869666.56 \\
    \hline
    \hline
    %
    \rowcolor{lightgray} 50 & 8192 & write sequential & 35.30 & 29.50 & 32114601.56 \\
    \hline
    \rowcolor{lightgray} 50 & 8192 & read sequential & 503.89 & 19.70 & 1024655425.65 \\
    \hline
    \rowcolor{lightgray} 50 & 8192 & write random & 33.69 & 28.20 & 36443371.02 \\
    \hline
    \rowcolor{lightgray} 50 & 8192 & read random & 504.57 & 18.20 & 1317877308.39 \\
    \hline
    \hline
    %
    50 & 16384 & write sequential & 35.05 & 29.75 & 35600309.01 \\
    \hline
    50 & 16384 & read sequential & 504.23 & 20.80 & 922037793.32 \\
    \hline
    50 & 16384 & write random & 36.36 & 26.90 & 34183453.23 \\
    \hline
    50 & 16384 & read random & 504.45 & 19.10 & 1338419510.38 \\
    \hline
    \hline
    %
    \rowcolor{lightgray} 100 & 1024 & write sequential & 31.12 & 29.15 & 33501250.50 \\
    \hline
    \rowcolor{lightgray} 100 & 1024 & read sequential & 504.33 & 20.60 & 1039245765.82 \\
    \hline
    \rowcolor{lightgray} 100 & 1024 & write random & 35.20 & 29.30 & 35365136.03 \\
    \hline
    \rowcolor{lightgray} 100 & 1024 & read random & 504.15 & 21.50 & 1149032035.95 \\
    \hline
    \hline
    %
    100 & 8192 & write sequential & 33.37 & 30.73 & 38188489.47 \\
    \hline
    100 & 8192 & read sequential & 495.67 & 21.30 & 1050523108.90 \\
    \hline
    100 & 8192 & write random & 38.09 & 30.97 & 38085067.97 \\
    \hline
    100 & 8192 & read random & 495.23 & 20.30 & 1350320079.31 \\
    \hline
    \hline
    %
    \rowcolor{lightgray} 100 & 16384 & write sequential & 30.89 & 30.90 & 35045035.83 \\
    \hline
    \rowcolor{lightgray} 100 & 16384 & read sequential & 496.12 & 21.90 & 1036803179.47 \\
    \hline
    \rowcolor{lightgray} 100 & 16384 & write random & 38.10 & 32.00 & 39262702.02 \\
    \hline
    \rowcolor{lightgray} 100 & 16384 & read random & 496.78 & 21.10 & 1179157844.37 \\
    \hline
    \hline
    \end{tabularx}

\section{Valutazione dei risultati}
Dai risultati dei test, è evidente che il costo prestazionale del TCFS è significativo. Tuttavia, è possibile mitigare questo costo attraverso l’implementazione della parallelizzazione del processo di cifratura, a condizione che le dimensioni dei blocchi siano sufficientemente grandi. Questo perché, nel caso di blocchi di dimensioni ridotte, l’overhead introdotto dalla parallelizzazione potrebbe effettivamente superare i vantaggi derivanti da tale approccio. Pertanto, è fondamentale bilanciare attentamente queste due considerazioni per ottimizzare le prestazioni del sistema. È interessante osservare che, in alcuni casi, il throughput è notevolmente elevato, superando di gran lunga le previsioni basate sui MB/s misurati. Questo fenomeno può essere attribuito con certezza a meccanismi di caching implementati da FUSE e/o NFS. Tali meccanismi, essendo inevitabili, possono influenzare significativamente le prestazioni del sistema. Inoltre si osserva che l’utilizzo della CPU aumenta con l’incremento della dimensione dei blocchi quando si utilizza TCFS. Parallelamente, anche la velocità di lettura e scrittura tende ad aumentare. Pertanto, la scelta della dimensione dei blocchi è un aspetto cruciale. Su sistemi con prestazioni inferiori, potrebbe essere preferibile optare per blocchi di dimensioni più piccole, anche se ciò potrebbe comportare una riduzione della velocità. Al contrario, su sistemi più performanti, si potrebbe considerare l’uso di blocchi di dimensioni maggiori.

\section{Possibili attacchi e contromisure}
Di seguito si presenta una lista dei possibili attacchi a TCFS:
\subsubsection{AES in modalità CBC (Cipher Block Chaining)}
\begin{itemize}
    \item \textbf{Padding Oracle Attack}: Questo attacco sfrutta il modo in cui alcuni sistemi rivelano se il padding di un messaggio crittografato è valido o meno\cite{AES-BasedTextEncryptionForNFC}. Un attaccante può usare queste informazioni per decifrare il messaggio byte per byte.\cite{cluster-basedPracticalKeyRecovery-differentialCryptanalysis}
    \item \textbf{IV Reuse}: Se l’IV viene riutilizzato con la stessa chiave, un attaccante può lanciare vari attacchi, tra cui attacchi di tipo chosen-plaintext. \cite{AES-CBC-ChosenPlaintextAttach}
\end{itemize}
\subsubsection{AES in modalità GCM (Galois/Counter Mode)}
\begin{itemize}
    \item \textbf{Forgery Attacks}: Se si riutilizza lo stesso IV per ogni crittografia eseguita con la stessa chiave, l’applicazione diventa vulnerabile agli attacchi di falsificazione. \cite{ForgeryAttacksOnGCM}
    \item \textbf{Ferguson and Saarinen Attacks}: Questi attacchi sfruttano le debolezze nella generazione del tag di autenticazione in GCM. \cite{CyclicAttackMACs}
\end{itemize}
\subsubsection{AES in modalità CTR (Counter Mode)}
\begin{itemize}
    \item \textbf{IV Reuse}: Come per la modalità CBC, il riutilizzo dell’IV può portare a vari attacchi, tra cui attacchi di tipo chosen-plaintext. \cite{ZeroDaysAttachDetection}
    \item \textbf{Malleability}: In modalità CTR, un attaccante può modificare i bit del testo cifrato senza conoscere la chiave, alterando così il testo in chiaro. \cite{ZeroDaysAttachDetection}
\end{itemize}
\subsubsection{AES senza IV}
L’uso di AES senza un IV può portare a vari problemi. Ad esempio, se si crittografa lo stesso messaggio più volte con la stessa chiave, il testo cifrato sarà lo stesso ogni volta. Questo può rivelare informazioni all’attaccante. \cite{DeepL-AES-KeyDecryption}
Inoltre, AES è resistente agli attacchi di tipo known-plaintext solo se si utilizza un IV diverso e casuale per ogni singolo messaggio. \cite{SideChannelAnalisys}

\section{Valutazioni personali}
A mio avviso, questa “proof of concept” ha superato le aspettative, dimostrando la sua efficacia. In un contesto di utilizzo reale, l’adozione di TCFS non sembra comportare una riduzione delle prestazioni tale da essere rilevante per l’utente. Questo suggerisce che TCFS può essere implementato con successo senza compromettere l’esperienza dell’utente.